\section{Certifying the Reader's Context}
\label{sec:method}

\subsection{Compile once, query each claim in constant time}

Figure~\ref{fig:compiler} gives the complete scan.
A stable source identifier breaks ties between witnesses.

\begin{figure}[t]
\centering
\fbox{\begin{minipage}{.92\columnwidth}
\small
\textbf{Input:} claims and accepted pairs $(d,c)$.\\
\textbf{Output:} boundaries and their source witnesses.
\begin{enumerate}
\item Set $p_c=r_c=\infty$ for every claim.
\item For each accepted pair with $d\ne c$:
\item\hspace{.4em} If $\ell_d>u_c$, minimize $p_c$ using $u_d$.
\item\hspace{.4em} If $u_d>\ell_c$, minimize $r_c$ using $\ell_d$.
\item Retain one supplying witness for each minimum.
\end{enumerate}
Steps 3--4 execute inside the pair scan.
\end{minipage}}
\caption{Linear compilation of the two support boundaries. The resulting records answer every query window.}
\label{fig:compiler}
\end{figure}

Let $n$ count modeled events, including explicit ends, and $m$ count accepted pairs.
Compilation takes $O(n+m)$ time and stores $O(n)$ certificate records.
A support decision takes $O(1)$ time per claim.
A complete mask takes $O(n)$ time.
Pair discovery remains the source adapter's responsibility.
The reference complete scan makes $n(n-1)$ gate calls.
An index can replace this scan when it preserves every accepted pair.

\subsection{An exact stability test}

Let $\pi_q$ be a complete relevance ranking for question $q$, including all tie breaks.
For a set $S$, write $\operatorname{Top}_k(S)$ for its first $k$ ranked claims.
The query's possible and guaranteed sets are $P(Q)$ and $H(Q)$.
A realized timeline $x$ supplies a support set $S_x(Q)$.
The reader's selected claim context is
\begin{equation}
 R_x(q)=\operatorname{Top}_k(S_x(Q)).
\label{eq:context}
\end{equation}
Filtering preserves the ranking and occurs before the budget.

\begin{theorem}[Exact context stability]
\label{thm:context-stability}
Every allowed timeline returns the same context exactly when
\begin{equation}
 \operatorname{Top}_k(P(Q))=\operatorname{Top}_k(H(Q)).
\label{eq:context-stability}
\end{equation}
Equivalently, every selected possible claim has guaranteed support.
\end{theorem}
\begin{proof}
Set $C=\operatorname{Top}_k(P(Q))$.
If $C\subseteq H(Q)$, every world retains $C$ and selects the same context.
Otherwise, choose $c\in C\setminus H(Q)$.
Some world supports $c$, while another does not.
At most $k-1$ possible claims outrank $c$.
Every world supporting $c$ therefore selects it.
The two worlds return different contexts.
For $k=0$, every context is empty.
\end{proof}

The test uses two filtered rankings without enumerating timelines.
Let $F$ contain the fixed fallback excerpts.
Apply the test to $P(Q)\cup F$ and $H(Q)\cup F$.
Their inclusion remains constant while modeled dates vary.
Claim selection precedes source rendering and any word-budget truncation.
We report rendered-context agreement separately.

\paragraph{Choose a stable context budget.}
The same scan finds the largest stable prefix.
Let $r$ be the position of the first unresolved claim in the possible-support ranking.
Every budget $k<r$ returns a fixed context, while every $k\ge r$ permits a context change.
If no unresolved claim exists, every budget is stable.
Each earlier modeled claim is guaranteed; earlier fallback excerpts stay eligible.
The claim at position $r$ supplies a counterexample for every budget that reaches it.
The system can therefore report a stable prefix alongside the unresolved evidence requiring date refinement.
This decision preserves the supplied relevance order.

\paragraph{Stop once the context is full.}
The implementation evaluates certificates lazily during selection.
It scans ranked candidates until $k$ possible excerpts fill the context.
It checks guaranteed support only for selected modeled claims.
Fixed fallback excerpts remain eligible without a temporal lookup.
If selection examines $j$ candidates, selection and certification take $O(j)$ time.
Other dates need no evaluation for this decision.

\subsection{Constructing and resolving an unstable context}
Let $U=\operatorname{Top}_k(P(Q)\cup F)\setminus(H(Q)\cup F)$.
These claims have unresolved support and enter the possible context.
Their status determines whether date uncertainty can change retrieval.

\begin{theorem}[Constructive instability and refinement]
\label{thm:refinement}
If $U\ne\varnothing$, we construct two complete date assignments with different contexts in $O(n)$ time.
Here, $n$ counts the event dates.
Every date refinement that stabilizes the context makes each claim in $U$ guaranteed or impossible.
\end{theorem}

Each $c\in U$ has fewer than $k$ possible claims above it.
Any timeline supporting $c$ therefore selects it.
We place every accepted replacement no later than $c$ or after a supporting query time.
A second timeline places $c$ after the query, or a direct witness by the query start.
Date refinement can only remove possible claims.
Thus a remaining possible pivot cannot leave the possible context through rank changes.
It must become guaranteed.
The assignments use only event bounds and the compact certificate.
Let $N$ count ranked excerpts, including fallback text.
Full context replay takes $O(n+m+N)$ time with the accepted pair list.
Appendix~\ref{app:refinement} gives the construction and proof.

The set $U$ identifies support decisions that every stabilizing refinement must settle.
New unresolved claims can enter after current pivots disappear.
We recompute the frontier after each refinement.

\subsection{A tractability boundary}
Possible support concerns one claim at a time.
Possible context membership also depends on which higher-ranked claims coexist.
We ask whether a given claim enters the context under \emph{any} allowed timeline.
\begin{theorem}[Possible context membership]
\label{thm:hardness}
When $k$ is part of the input, deciding possible context membership is NP-hard.
\end{theorem}
This holds even for a point query and a bipartite gate with one edge layer.
\begin{proof}
Take a Set Cover instance with elements $E$, sets $S_1,\ldots,S_s$, and budget $1\le B\le s$~\citep{karp1972reducibility}.
Each set becomes a claim $v_j$ with start bounds $[1,3]$.
Create $B+1$ copies of each element, all starting at time 0.
Add a candidate $c$ starting at time 0.
A pair $v_j\to e$ targets every copy of each element in $S_j$.
No other pairs are accepted.
Rank every set claim and element copy above $c$.
Use point query 2 and context budget $B+1$.

For a cover of at most $B$ sets, start its set claims at time 1.
Start the others at time 3.
Every element copy retires before the query.
At most $B$ higher-ranked set claims remain active, so $c$ enters the context.
Conversely, a context containing $c$ has at most $B$ active higher-ranked claims.
Its active set claims must cover every element.
An uncovered element would leave $B+1$ active copies above $c$ and exclude it.
Thus these claims define a cover of size at most $B$.
The construction is polynomial because $B\le s$.
\end{proof}
Appendix~\ref{app:refinement} also gives a reduction where every start precedes or equals the query time.
Our stability test instead needs one unresolved claim inside the possible prefix.
Its rank guarantees selection whenever it has support.
The test needs only individual support, even when shared events couple claim eligibility.
This makes exact stability and its witnessing timelines tractable without computing every attainable context member.


\subsection{What the witnesses explain}

Each certificate stores a target's bounds, both minima, and their supplying source identifiers.
An exclusion caused by retirement points to its direct replacement witness.
An exclusion caused by a late start points to the target's own date bounds.
This makes the temporal decision inspectable without reconstructing a cluster history.

The two boundaries answer different questions.
For $c=[0,2]$, take accepted witnesses $d=[1,4]$ and $e=[3,5]$.
Witness $d$ supplies $r_c=1$ because it can follow an early start of $c$.
Witness $e$ supplies $p_c=5$ because it must follow every start of $c$.
Keeping only $d$ incorrectly allows $c$ at time 5.
Keeping only $e$ incorrectly guarantees $c$ at time 2.
Neither one-witness subset reproduces every query decision.

The same direct structure limits update propagation.
Adding events whose earliest dates exceed the query end preserves that query's earlier support decisions.
Later ingestion can still correct a historical date when its event bounds reach the queried period.
These properties depend on direct temporal support, rather than document arrival order.
